\documentclass[11pt]{article}

\usepackage[margin=1in]{geometry}
\usepackage{amsmath,amssymb,amsthm,mathtools}
\usepackage{booktabs}
\usepackage{enumitem}
\usepackage{hyperref}
\usepackage{xcolor}

\hypersetup{
    colorlinks=true,
    linkcolor=blue,
    citecolor=blue,
    urlcolor=blue
}

\newtheorem{theorem}{Theorem}
\newtheorem{lemma}{Lemma}
\newtheorem{proposition}{Proposition}
\newtheorem{definition}{Definition}
\newtheorem{remark}{Remark}
\newtheorem{question}{Question}

\newcommand{\Fp}{\mathbb{F}_p}
\newcommand{\Fpx}{\mathbb{F}_p^\times}
\newcommand{\PhiPot}{\Phi}

\title{A Certified Finite-Check Package for Complement Cases in Erd\H{o}s Problem 475}
\author{Terence Fisher}
\date{\today}

\begin{document}

\maketitle

\begin{abstract}
Erd\H{o}s Problem 475 asks whether, for every prime \(p\) and every finite set
\(A\subseteq \mathbb F_p\setminus\{0\}\), there is a rearrangement
\(A=\{a_1,\ldots,a_t\}\) such that the partial sums
\[
a_1,\quad a_1+a_2,\quad \ldots,\quad a_1+\cdots+a_t
\]
are pairwise distinct modulo \(p\).

This document gives a certified finite-domain proof for the complement cases
\[
p=29,\quad |A|\in\{21,22,23,24,25\},
\]
and
\[
p=31,\quad |A|\in\{24,25,26,27\}.
\]
Equivalently, if
\[
B=\mathbb F_p^\times\setminus A,
\]
then the certified complement sizes are
\[
p=29,\quad 3\le |B|\le 7,
\qquad
p=31,\quad 3\le |B|\le 6.
\]

The proof is computer-assisted but independently checkable. The verifier reconstructs each trace state, recomputes the collision potential
\[
\Phi=(E,P,L),
\]
applies the certified move or path, and checks strict lexicographic descent. A separate coverage audit confirms that all canonical multiplicative-scaling representatives in the audited complement domain are present.
\end{abstract}

\section{Status and Scope}

This paper proves a certified finite-domain theorem for the complement cases listed above. It should be read as a completion of a finite check, not as a proof of all analytic reductions leading to that finite check.

The external reduction question is the following.

\begin{quote}
After applying the known analytic results for Erd\H{o}s Problem 475, are the only remaining complement cases
\[
p=29,\quad 3\le |B|\le7,
\qquad
p=31,\quad 3\le |B|\le6,
\]
up to nonzero multiplicative scaling?
\end{quote}

Under this reduction, the present certificate completes the remaining finite verification.

\section{Problem Setup}

Let \(p\) be prime and let
\[
A\subseteq \Fpx.
\]
An ordering of \(A\) is a sequence
\[
R=(r_1,\ldots,r_t)
\]
whose entries are exactly the elements of \(A\).

\begin{definition}[Partial sums]
For an ordering \(R=(r_1,\ldots,r_t)\), define
\[
S_i(R)=\sum_{j=1}^{i}r_j \pmod p,
\qquad
1\le i\le t.
\]
\end{definition}

\begin{definition}[Valid ordering]
An ordering \(R\) is valid if
\[
S_1(R),S_2(R),\ldots,S_t(R)
\]
are pairwise distinct in \(\Fp\).
\end{definition}

The computational traces are written in terms of a complement set
\[
B=\Fpx\setminus A.
\]
The trace field \(Q_p\) is a permutation of \(\Fpx\), and the trace field \texttt{initial\_order} satisfies
\[
\texttt{initial\_order}=Q_p\setminus B.
\]
Therefore, the ordering being repaired is an ordering of the original Erd\H{o}s set \(A\), not of the small complement \(B\).

\section{Trace Semantics}

The trace-semantics verifier checks every row of the trace files and confirms the following facts.

\begin{enumerate}[label=(\roman*)]
\item \(Q_p\) is a permutation of \(\Fpx=\{1,\ldots,p-1\}\).
\item \(B\subseteq Q_p\).
\item \(\texttt{initial\_order}=Q_p\setminus B\).
\item \(\texttt{final\_order}\) is a permutation of \(\texttt{initial\_order}\).
\item \(\texttt{final\_partial\_sums}\) are the partial sums of \(\texttt{final\_order}\).
\item \(\texttt{final\_partial\_sums}\) are pairwise distinct.
\end{enumerate}

Thus each trace row explicitly represents a successful ordering of
\[
A=\Fpx\setminus B.
\]

The audited complement-to-original-size conversion is
\[
\begin{array}{c|c|c}
p & |B| & |A|=p-1-|B|\\
\hline
29 & 3,4,5,6,7 & 25,24,23,22,21\\
31 & 3,4,5,6 & 27,26,25,24
\end{array}
\]

This places the certificate in the near-full-set regime, not in the small-\(|A|\) regime.

\section{Collision Potential}

For \(x\in\Fp\), define
\[
m_x(R)=|\{i:S_i(R)=x\}|.
\]

\begin{definition}[Excess collision count]
Define
\[
E(R)=\sum_{x\in\Fp}\max(m_x(R)-1,0).
\]
Then \(E(R)=0\) if and only if \(R\) is valid.
\end{definition}

\begin{definition}[Collision pair count]
Define
\[
P(R)=\sum_{x\in\Fp}\binom{m_x(R)}{2}.
\]
\end{definition}

\begin{definition}[Collision span weight]
Define
\[
L(R)=
\sum_{\substack{1\le i<j\le t\\S_i(R)=S_j(R)}}(j-i).
\]
\end{definition}

\begin{definition}[Lexicographic potential]
Define
\[
\PhiPot(R)=(E(R),P(R),L(R)),
\]
ordered lexicographically.
\end{definition}

The proof strategy is to rule out a minimal invalid ordering by showing that every invalid certified state admits either strict descent in \(\PhiPot\), or a bounded local escape path ending in strict descent.

\section{Scaling Symmetry}

\begin{lemma}[Multiplicative scaling invariance]
Let \(a\in\Fpx\). For an ordering
\[
R=(r_1,\ldots,r_t),
\]
define
\[
aR=(ar_1,\ldots,ar_t).
\]
Then
\[
E(aR)=E(R),\qquad P(aR)=P(R),\qquad L(aR)=L(R).
\]
Therefore
\[
\PhiPot(aR)=\PhiPot(R).
\]
\end{lemma}

\begin{proof}
The partial sums of \(aR\) are
\[
S_i(aR)=\sum_{j=1}^{i}ar_j=aS_i(R).
\]
Since \(a\ne0\), multiplication by \(a\) is a bijection of \(\Fp\). Hence
\[
S_i(R)=S_j(R)
\quad\Longleftrightarrow\quad
S_i(aR)=S_j(aR).
\]
Thus the equality pattern among partial sums is unchanged. The quantities \(E\), \(P\), and \(L\) depend only on this equality pattern and the index gaps \(j-i\). Hence they are preserved.
\end{proof}

Consequently, it is sufficient to check one canonical representative of each orbit
\[
B\sim aB,\qquad a\in\Fpx.
\]

\section{Local Block Moves}

A local move is a contiguous block move. If
\[
R=X\,U\,Y\,Z,
\]
where \(U\) is a contiguous block, then the move sends
\[
X\,U\,Y\,Z\longmapsto X\,Y\,U\,Z.
\]

The scripts encode such a move by
\[
(s,\ell,t),
\]
where \(s\) is the starting index of \(U\), \(\ell=|U|\), and \(t\) is the insertion index after removing \(U\).

The certificate uses two locality radii:
\[
h=4
\]
for the \(E\ge2\) descent branches, and
\[
h=6
\]
for the one-collision escape branch.

\section{Branch Decomposition}

Every invalid state has
\[
E(R)>0.
\]
The certificate decomposes invalid states into three branches.

\subsection{Non-atomic \(E\ge2\)}

The non-atomic branch consists of states satisfying
\[
E(R)\ge2
\]
and not satisfying
\[
P(R)=2,\qquad \max_xm_x(R)=2.
\]

In the verified certificate, these states admit halo-4 local descent moves anchored at maximum-multiplicity collision structure. In the full \(p=29\) file, every non-atomic state is resolved by a max-multiplicity halo-4 singleton move. In the full \(p=31\) file, every non-atomic state is resolved, with only one requiring a block of length two.

\subsection{Atomic \(E\ge2\)}

The atomic branch is
\[
P(R)=2,\qquad \max_xm_x(R)=2.
\]
It has exactly two collision pairs and no residue appearing more than twice.

This branch is certified by a finite table of relative halo-4 moves with
\[
|U|\le2.
\]
The compressed atomic macro-certificate contains \(54\) relative move rules covering all verified atomic instances.

\subsection{One-collision branch}

The one-collision branch is
\[
E(R)=P(R)=1.
\]

Direct monotone descent is not always available here. The certificate allows a bounded path
\[
R=R_0\to R_1\to\cdots\to R_t
\]
with
\[
t\le3.
\]
The intermediate states satisfy
\[
E(R_i)=P(R_i)=1,\qquad i<t,
\]
and the final state satisfies
\[
\PhiPot(R_t)<\PhiPot(R_0).
\]
All moves in this path are halo-6 moves with
\[
|U|\le5.
\]

\section{Certified Branch Theorem}

\begin{theorem}[Certified descent and escape theorem]
For every invalid certified state \(R\) in the audited complement domain, one of the following holds.

\begin{enumerate}[label=(\roman*)]
\item If
\[
E(R)\ge2,\qquad (P(R),\max_xm_x(R))\ne(2,2),
\]
then there exists a certified halo-4 local move \(R\to R'\) such that
\[
\PhiPot(R')<\PhiPot(R).
\]

\item If
\[
P(R)=2,\qquad \max_xm_x(R)=2,
\]
then one of the certified atomic halo-4 moves with \(|U|\le2\) gives
\[
\PhiPot(R')<\PhiPot(R).
\]

\item If
\[
E(R)=P(R)=1,
\]
then there exists a certified path
\[
R=R_0\to R_1\to\cdots\to R_t,\qquad t\le3,
\]
such that
\[
E(R_i)=P(R_i)=1,\qquad i<t,
\]
and
\[
\PhiPot(R_t)<\PhiPot(R_0).
\]
Each move in this path is halo-6 local and moves a block of length at most five.
\end{enumerate}
\end{theorem}

\begin{proof}
This theorem is certified by the supplied CSV files and independently checked by the verifier.

For each certified row, the verifier:
\begin{enumerate}
\item reconstructs the referenced ordering \(R\) from the original JSONL trace;
\item recomputes \(E(R)\), \(P(R)\), \(L(R)\), and \(\PhiPot(R)\);
\item checks the branch classification;
\item applies the claimed move or path;
\item recomputes the final potential;
\item checks the strict inequality
\[
\PhiPot(R_{\mathrm{final}})<\PhiPot(R).
\]
\end{enumerate}

For the one-collision branch, it additionally verifies that intermediate states remain one-collision:
\[
E=P=1.
\]
The independent verifier reported
\[
198{,}631
\]
successful checks and
\[
0
\]
failures.
\end{proof}

\section{Certificate Statistics}

\[
\begin{array}{lrr}
\toprule
\text{Certificate family} & \text{Verified rows} & \text{Failures}\\
\midrule
\text{Non-atomic }E\ge2\text{ descent} & 101{,}385 & 0\\
\text{One-collision escape} & 39{,}812 & 0\\
\text{Atomic instance candidates} & 28{,}717 & 0\\
\text{Atomic signature certificates} & 28{,}717 & 0\\
\midrule
\text{Total} & 198{,}631 & 0\\
\bottomrule
\end{array}
\]

The non-atomic family decomposes as:
\[
65{,}272
\]
states for \(p=29\), and
\[
36{,}113
\]
states for \(p=31\).

The one-collision family decomposes as:
\[
29{,}854
\]
states for \(p=29\), and
\[
9{,}958
\]
states for \(p=31\).

\section{Coverage of Canonical Representatives}

\begin{theorem}[Coverage theorem]
The trace files contain every canonical multiplicative-scaling representative in the audited complement domain:
\[
p=29,\quad 3\le |B|\le7,
\]
and
\[
p=31,\quad 3\le |B|\le6.
\]
\end{theorem}

\begin{proof}
The coverage audit enumerates all canonical representatives of the action
\[
B\mapsto aB,\qquad a\in\Fpx,
\]
and compares them with the representatives present in the trace files.

For every audited pair \((p,|B|)\), the canonical-scaling audit reports:
\[
\text{coverage}=100\%,\qquad \text{missing}=0,\qquad \text{extra}=0.
\]
By multiplicative scaling invariance, this proves coverage of the audited complement domain modulo the valid symmetry \(B\sim aB\).
\end{proof}

\section{Minimal-Obstruction Proof}

\begin{theorem}[Certified finite complement theorem]
Let
\[
A\subseteq \Fpx
\]
be in one of the following cases:
\[
p=29,\quad |A|\in\{21,22,23,24,25\},
\]
or
\[
p=31,\quad |A|\in\{24,25,26,27\}.
\]
Equivalently, let
\[
B=\Fpx\setminus A
\]
satisfy
\[
p=29,\quad 3\le |B|\le7,
\]
or
\[
p=31,\quad 3\le |B|\le6.
\]
Then, for every canonical multiplicative-scaling representative \(B\), there exists an ordering \(R\) of \(A\) such that
\[
S_1(R),\ldots,S_{|A|}(R)
\]
are pairwise distinct.
\end{theorem}

\begin{proof}
Suppose, for contradiction, that some audited complement case has no valid ordering. Among all certified invalid states for this case, choose \(R\) with minimal \(\PhiPot(R)\) in lexicographic order.

Since \(R\) is invalid,
\[
E(R)>0.
\]

If
\[
E(R)\ge2,
\]
then \(R\) lies either in the non-atomic branch or in the atomic branch. In both cases, the certified branch theorem gives a state \(R'\) satisfying
\[
\PhiPot(R')<\PhiPot(R),
\]
contradicting the minimality of \(R\).

It remains to consider
\[
E(R)=1.
\]
Then exactly one residue occurs twice and all others occur at most once. Therefore
\[
P(R)=1.
\]
Thus \(R\) lies in the one-collision branch. The certified branch theorem gives a path
\[
R=R_0\to R_1\to\cdots\to R_t
\]
with
\[
t\le3
\]
and
\[
\PhiPot(R_t)<\PhiPot(R_0).
\]
This again contradicts the minimality of \(R\).

Therefore no invalid minimal obstruction exists in the audited complement domain. Hence every audited \(A\) has a valid ordering.
\end{proof}

\section{Reduction Status Relative to the Published Problem}

The public Erd\H{o}s Problems database states Problem \#475 in the following form: for a prime \(p\), every finite set
\[
A\subseteq \mathbb F_p\setminus\{0\}
\]
should admit a rearrangement
\[
A=\{a_1,\ldots,a_t\}
\]
such that the partial sums
\[
a_1,\quad a_1+a_2,\quad \ldots,\quad a_1+\cdots+a_t
\]
are pairwise distinct modulo \(p\).

The same source currently records the problem as ``resolved up to a finite check.'' It lists the following known ranges:
\[
t\le 12,
\]
\[
p-3\le t\le p-1,
\]
and all sufficiently large primes, obtained by combining known small, medium, large, and very-large set-size regimes.

The present certificate verifies the finite complement cases
\[
p=29,\quad |A|=21,22,23,24,25,
\]
and
\[
p=31,\quad |A|=24,25,26,27.
\]
Equivalently, with
\[
B=\mathbb F_p^\times\setminus A,
\]
the certified complement sizes are
\[
p=29,\quad |B|=3,4,5,6,7,
\]
and
\[
p=31,\quad |B|=3,4,5,6.
\]

These cases lie immediately below the already-known very-large range
\[
p-3\le |A|\le p-1.
\]
For \(p=29\), that known range begins at \(|A|=26\). For \(p=31\), it begins at \(|A|=28\).

Thus the certificate is not merely a recheck of the known small-set range \(t\le12\). It verifies large original sets \(A\) with small complements.

The remaining reduction question is the following.

\begin{quote}
Do the published analytic reductions leave exactly the complement cases
\[
p=29,\quad |B|=3,\ldots,7,
\qquad
p=31,\quad |B|=3,\ldots,6
\]
as the finite check?
\end{quote}

The certificate package proves these finite cases. It does not independently prove that no other finite cases remain after the analytic reductions.

\section{Connection to Erd\H{o}s Problem 475}

The certified theorem above proves the finite complement cases explicitly represented by the traces. To upgrade this to a full proof of Erd\H{o}s Problem 475, one must combine it with the external analytic reduction showing that all other cases are already covered by known results.

The current package therefore establishes the computational finite-check component:
\[
p=29,\quad |A|=21,\ldots,25,
\]
and
\[
p=31,\quad |A|=24,\ldots,27.
\]

If these are exactly the finite residue cases left by the known analytic reductions, then this certificate completes the remaining finite check.

\section{Reproducibility Commands}

The trace-semantics verification command is:
\begin{verbatim}
python scripts/verify_erdos475_trace_semantics.py \
  traces/p29_r3_to_r7_repair_traces_strict.jsonl \
  traces/p31_r3_to_r6_repair_traces_strict.jsonl
\end{verbatim}

The expected result is:
\begin{verbatim}
VERDICT: PASS
\end{verbatim}

The certificate-verification command is:
\begin{verbatim}
python scripts/verify_erdos475_certificates.py \
  --trace-files \
    traces/p29_r3_to_r7_repair_traces_strict.jsonl \
    traces/p31_r3_to_r6_repair_traces_strict.jsonl \
  --nonatomic-csv \
    certificates/p29_nonatomic_descent_full.csv \
    certificates/p31_nonatomic_descent_full.csv \
  --onecollision-csv \
    certificates/p29_one_collision_deep_full.csv \
    certificates/p31_one_collision_deep_full.csv \
  --atomic-instances certificates/atomic_local_cert_instances.csv \
  --atomic-certs certificates/atomic_local_certs.csv \
  --require-onecollision-intermediates \
  --strict-csv
\end{verbatim}

The expected result is:
\begin{verbatim}
TOTAL ok=198631 fail=0
VERDICT: PASS
\end{verbatim}

The coverage-audit command is:
\begin{verbatim}
python scripts/audit_erdos475_trace_coverage.py \
  traces/p29_r3_to_r7_repair_traces_strict.jsonl \
  traces/p31_r3_to_r6_repair_traces_strict.jsonl \
  --summary-csv certificates/trace_coverage_summary.csv
\end{verbatim}

The expected canonical-scaling result is:
\begin{verbatim}
coverage=100.00%
missing=0
extra=0
\end{verbatim}

\end{document}
